\section{Заключение}
\label{sec:Chapter5}

В работе разработана ультразвуковая локальная навигационная система для
определения координат привязного БПЛА относительно наземной платформы в условиях
недоступности глобальных навигационных спутниковых систем. Получены следующие
основные результаты.

\begin{enumerate}
  \item Сформулирована постановка задачи и требования к системе
        (раздел~\ref{sec:Chapter1}); проведён сравнительный анализ существующих
        подходов к локализации БПЛА, обосновавший выбор архитектуры~--- ультразвуковая
        дальнометрия с синхронизацией по UWB-радиоканалу (раздел~\ref{sec:Chapter2}).

  \item Построена математическая модель системы (раздел~\ref{sec:Chapter3}):
        \begin{itemize}
          \item оценена максимальная дальность работы (${\approx}35$~м для текущего
                излучателя на частоте 25~кГц; целевая дальность~--- $40$--$50$~м);
          \item проанализирована погрешность определения скорости звука во влажном
                воздухе и показано, что при использовании недорогих датчиков она не
                превышает ${\sim}0{,}13\%$;
          \item получено выражение для влияния геометрии якорей (GDOP) на ошибку
                позиционирования;
          \item проведён сравнительный анализ аналитического метода трилатерации и
                метода Гаусса--Ньютона; показано, что для конфигурации из четырёх
                якорей предпочтителен вычислительно простой аналитический метод.
        \end{itemize}

  \item Разработана и собрана аппаратная часть (раздел~\ref{sec:Chapter4}):
        передающий узел БПЛА с аппаратной синхронизацией ультразвукового и
        UWB-сигналов по EXTTXE, тракт обработки сигнала якоря (предусилитель,
        полосовой усилитель, детектор огибающей и компаратор), двухимпульсный метод
        самокалибровки времени прихода, устраняющий зависящее от дальности
        систематическое смещение компаратора, а также изолированная система питания
        с защитой от наводок. Описан протокол UWB-синхронизации узлов.

  \item Экспериментально подтверждена работоспособность приёмного тракта: уверенный
        приём достигнут на дистанции ${\approx}34$~м (ограничена размерами площадки
        испытаний, а не системой), сигнал разрешается в пару импульсов, а время
        нарастания фронта огибающей составило ${\approx}0{,}2$~мс. Расчётная погрешность
        позиционирования при достигнутых параметрах измерений~--- порядка $5$~см.
\end{enumerate}

Полученные результаты показывают, что поставленная задача решена: предложенная
система удовлетворяет сформулированным требованиям по дальности, точности и
условиям эксплуатации и доведена до уровня работающего прототипа. Дальнейшая работа
связана с полевыми испытаниями полной системы из четырёх якорей на открытом
пространстве, интеграцией фильтра оценивания траектории и проверкой работы вблизи
силового оборудования платформы.
